% 原创TikZ重绘；层次布局参考Bishop与Bishop（2024）第12章。
\begin{tikzpicture}[x=1mm,y=1mm,font=\figurefont\footnotesize,
 every node/.style={text=NNInk},
 stage/.style={nn operator,minimum height=9mm,inner sep=2pt},
 norm/.style={stage,minimum height=7mm},
 wire/.style={nn flow,draw=NNWire,rounded corners=1.5mm},
 aux/.style={draw=NNLine,line width=.85pt,densely dashed},
 sum/.style={nn pointwise,minimum size=5mm}]

\node[nn heading] (posttitle) at (39,137) {Post-LN};
\node[] (postin) at (39,17) {$\bm H^{(l-1)}$};
\node[stage,minimum width=39mm] (post0) at (39,34) {双向多头注意力};
\draw[wire] (postin.north)--(post0.south);
\node[sum] (post1) at (39,49) {$+$};
\draw[wire] (post0.north)--(post1.south);
\node[norm,minimum width=39mm] (post2) at (39,62) {LN};
\draw[wire] (post1.north)--(post2.south);
\node[stage,minimum width=39mm] (post3) at (39,82) {逐位置FFN};
\draw[wire] (post2.north)--(post3.south);
\node[sum] (post4) at (39,97) {$+$};
\draw[wire] (post3.north)--(post4.south);
\node[norm,minimum width=39mm] (post5) at (39,110) {LN};
\draw[wire] (post4.north)--(post5.south);
\node[] (postout) at (39,124) {$\bm H^{(l)}$};
\draw[wire] (post5.north)--(postout.south);
\draw[wire] (39,24)--(63.5,24)--(63.5,49)--(post1.east);
\fill[NNWire] (39,24) circle (.55);
\draw[wire] (39,70)--(63.5,70)--(63.5,97)--(post4.east);
\fill[NNWire] (39,70) circle (.55);
\node[inner sep=0pt,minimum width=12mm,minimum height=9mm] (postmatrix) at (39,5) {};
\filldraw[draw=NNInput,fill=NNInput!18,line width=.5pt] (33.0,0.5) rectangle ++(3.0,3);
\filldraw[draw=NNInput,fill=NNInput!18,line width=.5pt] (36.0,0.5) rectangle ++(3.0,3);
\filldraw[draw=NNInput,fill=NNInput!18,line width=.5pt] (39.0,0.5) rectangle ++(3.0,3);
\filldraw[draw=NNInput,fill=NNInput!18,line width=.5pt] (42.0,0.5) rectangle ++(3.0,3);
\filldraw[draw=NNInput,fill=NNInput!18,line width=.5pt] (33.0,3.5) rectangle ++(3.0,3);
\filldraw[draw=NNInput,fill=NNInput!18,line width=.5pt] (36.0,3.5) rectangle ++(3.0,3);
\filldraw[draw=NNInput,fill=NNInput!18,line width=.5pt] (39.0,3.5) rectangle ++(3.0,3);
\filldraw[draw=NNInput,fill=NNInput!18,line width=.5pt] (42.0,3.5) rectangle ++(3.0,3);
\filldraw[draw=NNInput,fill=NNInput!18,line width=.5pt] (33.0,6.5) rectangle ++(3.0,3);
\filldraw[draw=NNInput,fill=NNInput!18,line width=.5pt] (36.0,6.5) rectangle ++(3.0,3);
\filldraw[draw=NNInput,fill=NNInput!18,line width=.5pt] (39.0,6.5) rectangle ++(3.0,3);
\filldraw[draw=NNInput,fill=NNInput!18,line width=.5pt] (42.0,6.5) rectangle ++(3.0,3);
\node[anchor=east] (postxlabel) at (29,5) {输入矩阵};
\draw[wire] (postmatrix.north)--(postin.south);
\node[nn heading] (pretitle) at (119,137) {Pre-LN};
\node[] (prein) at (119,17) {$\bm H^{(l-1)}$};
\node[norm,minimum width=39mm] (pre0) at (119,31) {LN};
\draw[wire] (prein.north)--(pre0.south);
\node[stage,minimum width=39mm] (pre1) at (119,45) {双向多头注意力};
\draw[wire] (pre0.north)--(pre1.south);
\node[sum] (pre2) at (119,59) {$+$};
\draw[wire] (pre1.north)--(pre2.south);
\node[norm,minimum width=39mm] (pre3) at (119,73) {LN};
\draw[wire] (pre2.north)--(pre3.south);
\node[stage,minimum width=39mm] (pre4) at (119,87) {逐位置FFN};
\draw[wire] (pre3.north)--(pre4.south);
\node[sum] (pre5) at (119,101) {$+$};
\draw[wire] (pre4.north)--(pre5.south);
\node[] (preout) at (119,124) {$\bm H^{(l)}$};
\draw[wire] (pre5.north)--(preout.south);
\draw[wire] (119,24)--(143.5,24)--(143.5,59)--(pre2.east);
\fill[NNWire] (119,24) circle (.55);
\draw[wire] (119,66)--(143.5,66)--(143.5,101)--(pre5.east);
\fill[NNWire] (119,66) circle (.55);
\node[inner sep=0pt,minimum width=12mm,minimum height=9mm] (prematrix) at (119,5) {};
\filldraw[draw=NNInput,fill=NNInput!18,line width=.5pt] (113.0,0.5) rectangle ++(3.0,3);
\filldraw[draw=NNInput,fill=NNInput!18,line width=.5pt] (116.0,0.5) rectangle ++(3.0,3);
\filldraw[draw=NNInput,fill=NNInput!18,line width=.5pt] (119.0,0.5) rectangle ++(3.0,3);
\filldraw[draw=NNInput,fill=NNInput!18,line width=.5pt] (122.0,0.5) rectangle ++(3.0,3);
\filldraw[draw=NNInput,fill=NNInput!18,line width=.5pt] (113.0,3.5) rectangle ++(3.0,3);
\filldraw[draw=NNInput,fill=NNInput!18,line width=.5pt] (116.0,3.5) rectangle ++(3.0,3);
\filldraw[draw=NNInput,fill=NNInput!18,line width=.5pt] (119.0,3.5) rectangle ++(3.0,3);
\filldraw[draw=NNInput,fill=NNInput!18,line width=.5pt] (122.0,3.5) rectangle ++(3.0,3);
\filldraw[draw=NNInput,fill=NNInput!18,line width=.5pt] (113.0,6.5) rectangle ++(3.0,3);
\filldraw[draw=NNInput,fill=NNInput!18,line width=.5pt] (116.0,6.5) rectangle ++(3.0,3);
\filldraw[draw=NNInput,fill=NNInput!18,line width=.5pt] (119.0,6.5) rectangle ++(3.0,3);
\filldraw[draw=NNInput,fill=NNInput!18,line width=.5pt] (122.0,6.5) rectangle ++(3.0,3);
\node[anchor=east] (prexlabel) at (109,5) {输入矩阵};
\draw[wire] (prematrix.north)--(prein.south);
\draw[nn gradient,line width=1.35pt,rounded corners=1.2mm]
  (postout.east)--(69,124)--(69,14)--(postin.east);
\node[anchor=west,text=FigureInk,fill=none,inner sep=1pt,align=left]
  at (69,52) {$\partial L/\partial\bm x$\\仍经过LN};
\draw[nn gradient,line width=1.35pt,rounded corners=1.2mm]
  (preout.east)--(152,124)--(152,14)--(prein.east);
\node[anchor=west,text=FigureInk,fill=none,inner sep=1pt,align=left]
  at (152,52) {$\partial L/\partial\bm x$\\含恒等直通};
\end{tikzpicture}
